\section{Глава 1}

\todo[inline]{Придумать название для первой главы, корректно отражающую ее суть}

\subsection{Терминология предметной области исследования}

Ряд определений термина <<база данных>> достаточно обширен,
среди прочих он содержит несколько известных трактовок, однако
единой общепризнанной формулировки не существует. 
Так, в интерпретации Д. Кренке [1] \textit{база данных} --- это 
самодокументированное собрание интегрированных записей. 
В контексте определении базы данных как
<<собрания интегрированных записей>> она рассматривается в качестве
набора данных, связанных между собой некоторым образом. 
Свойство <<самодокументированности>> предполагает, что такой набор
данных содержит описание собственной структуры.  

\todo[inline]{Разобраться со списком источников и как правильно им пользоваться в latex}

\todo[inline, backgroundcolor=blue!20!white]{Ввести определения данных и информации, чтобы четко разграничить их}

\todo[inline, backgroundcolor=blue!20!white]{Определение модели данных}

\todo[inline, backgroundcolor=blue!20!white]{Определение метаданных}

\todo[inline, backgroundcolor=blue!20!white]{Определение понятий атрибут, запись, сущность, кортеж}

\todo[inline, backgroundcolor=blue!20!white]{Определение связи и отношения}

\todo[inline, backgroundcolor=blue!20!white]{Определение узла, ребра, свойства}

\todo[inline, backgroundcolor=blue!20!white]{Определение иерархии}

\todo[inline, backgroundcolor=blue!20!white]{Определение целостности данных}

\todo[inline, backgroundcolor=blue!20!white]{Определение семантической целостности}

\todo[inline, backgroundcolor=blue!20!white]{Определение транзакционности}

\todo[inline, backgroundcolor=blue!20!white]{Определение отображения моделей}

\todo[inline, backgroundcolor=blue!20!white]{Определение сериализации}

\subsection{Реляционная модель данных}

\subsection{Графовая модель данных}

\subsection{Сетевая модель данных}

\subsection{Иерархическая модель данных}